% =====================================================================
\section{Introduction}
\label{sec:intro}
% =====================================================================

\subsection{The Clay statement, exactly}

The Clay problem, as posed by Jaffe and Witten, is:

\begin{qbox}
Prove that for any compact simple gauge group $G$, a non-trivial quantum Yang--Mills theory
exists on $\mathbb{R}^4$ and has a mass gap $\Delta>0$. Existence includes establishing
axiomatic properties at least as strong as those cited in Streater--Wightman (1964) and
Osterwalder--Schrader (1973, 1975). [Jaffe--Witten, \emph{Quantum Yang--Mills Theory}, Clay
Mathematics Institute]
\end{qbox}

Two features of this statement are habitually under-weighted and drive everything that
follows. First, it demands a \textbf{uniform} gap: not merely that each finite lattice volume
has a positive spectral gap (trivially true by compactness), but that the gap survives
$L\to\infty$ and $a\to0$ with a volume- and cutoff-uniform lower bound, and that the resulting
theory clusters exponentially at a rate $m>0$. Second, it demands that the gap be a feature of
the \textbf{gauge-invariant, physical} theory --- not of a gauge-fixed chart, and not of a
perturbative expansion around a chosen background.

The campaign's operational form of the Clay conclusion, assembled and partially
machine-verified in \code{docs/conditional-theorem.md}, is:
\begin{align*}
&\exists\,\mu\ \text{on}\ \mathbb{R}^4,\ \text{gauge-invariant, satisfying OS0--OS4};\\
&\text{non-Gaussian},\quad g_{\mathrm{ren}}/\Lambda>0;\\
&\Delta/\Lambda\ \ge\ \min(\delta_W,\delta_S)>0;\\
&|S(f\,\tau_a g)-S(f)\,S(g)|\ \le\ C\,e^{-m\|a\|},\qquad
  m\ge c\,\Delta>0,\ \ \text{all gauge-invariant local } f,g;\\
&\text{unique vacuum};\quad \text{area law } \sigma=c_\sigma\Lambda^2>0;\quad
  \text{all constants group-uniform in } (h^{\vee},\dim G,\dim\mathfrak{g}).
\end{align*}

The transfer-matrix form of the same statement is the one used throughout this paper: for the
lattice transfer matrix $T_a=e^{-H_a}$ on the gauge-invariant Hilbert space $\mathcal{H}_a$,
with vacuum $\Omega$,
\begin{equation}
\exists\,\delta>0:\qquad
\inf\operatorname{spec}\!\bigl(H_a\big|_{\Omega^\perp}\bigr)\ \ge\ \delta\Lambda
\qquad\text{for all sufficiently small } a .
\label{eq:clay-transfer}
\end{equation}

\subsection{The lock-core rewriting}

The campaign works with a sharper, channel-specific rewriting of the gap which, while not
literally equivalent to the Clay statement, is the correct object to attack. Let $\nu$ be the
Källén--Lehmann spectral measure of the gauge-invariant local operator
$O=\operatorname{tr}F^2$ (more precisely $O_S=\tfrac14\operatorname{tr}F_{\mu\nu}F^{\mu\nu}$,
the scalar density; the normalization is immaterial). Reflection positivity gives $\nu\ge0$
and $\operatorname{supp}\nu\subset[m_1^2,\infty)$, where $m_1$ is the lightest mass in the
$0^{++}$ channel. The \textbf{atom} is
\begin{equation}
f_1^2 \;:=\; \nu(\{m_1^2\}) \;=\; \bigl|\langle 0|\operatorname{tr}F^2|\mathbb{G}\rangle\bigr|^2 ,
\label{eq:atom}
\end{equation}
the squared overlap of the vacuum with the lightest $0^{++}$ glueball $\mathbb{G}$. The \textbf{lock core}
is: $f_1^2>0$ and the atom is isolated ($\Delta>0$ above it). The campaign summarises this as
\begin{equation}
\text{Clay}\iff\neg(\mathrm{iii}),\qquad
(\mathrm{iii})=\text{``the YM measure is a gapless interacting 4D CFT''.}
\label{eq:lockcore}
\end{equation}

This rewriting is useful because it names the exact analytic object (an atom weight in a
positive measure at the bottom of a channel) and the exact analytic property it has: $f_1^2$ is
\textbf{non-diagonal} (it is an off-diagonal matrix element squared),
\textbf{upper-semicontinuous} in the coupling ($\nu(\{x\})$ is usc), and requires a
\textbf{second scale} $E_2$ to certify isolation. Each of these properties will matter: usc
means continuity can only make an atom \emph{fall}, never rise; the need for $E_2$ is the
reason every bottom-functional tool is blind.

\textbf{Honesty checkpoint.} ``Clay $\iff$ atom'' is \emph{not} a theorem of equivalence. The
campaign's own adversarial audit (\code{research/unconditional/REDUCTION-AUDIT.md}) establishes
that the correct logical shape is a \textbf{diamond}, not a chain (Figure~\ref{fig:diamond}): a
gapped theory need not have an atom at all (the massive free scalar with $O={:}\varphi^2{:}$
has a gap and $f_1^2=0$), so $\text{gap}\Rightarrow\text{atom}$ is false; the true direction is
\emph{isolated atom at $\inf\operatorname{supp}\nu$ $\Rightarrow$ exponential clustering in that
channel}, conditional on that channel carrying the lightest state. This paper therefore treats
``the atom $f_1^2>0$'' as the \textbf{target object}, not as a synonym for Clay. The gap
between the two is itself a named missing step --- the ordering input that the lightest state
lies in $0^{++}$ (the corpus's \code{IN-11}, a physical input, not derivable from symmetry).

\figreduction

\subsection{The two halves}

The campaign splits Clay into two halves and a channel:

\begin{itemize}
\item \textbf{G1 (existence).} 4D continuum existence with OS axioms. The fixed-volume
  construction is essentially complete; the one genuinely open UV input is a \textbf{uniform
  renormalization constant} $Z_a$, equivalently a cutoff-uniform bound on the force Hessian
  (the node called L3). L3 is purely ultraviolet and \emph{orthogonal} to the gap --- proving it
  would make fixed-volume existence unconditional but would not help G2.
\item \textbf{G2 (gap).} The proof or transport of $f_1^2>0$ (equivalently of the coarse
  curvature $G''(0)>0$ at $k\to0$). This is where the problem lives.
\item \textbf{The physical channel.} Identifying the lightest state as $0^{++}$ and showing it
  does not dissolve (the no-dissolution theorem). This half is, conditionally, complete.
\end{itemize}

The single deepest structural finding of the campaign is that \textbf{G1's L3 and G2's H1 are
one species} (Figure~\ref{fig:scales}): both are a one-sided \emph{upper bound on a variance}.
L3 is the UV force-variance bound; H1 is the IR conditioned-variance bound
$\operatorname{Var}(\xi_{k\to0})\le1/\lambda$ (curvature form
$\operatorname{Var}_\mu(\langle v,\nabla S\rangle|0)\le E_\mu[H_{vv}|0]+W_{\mathrm{full}}-\delta$).
The shapes are
identical because $\operatorname{Var}\ge0$ can only supply a lower edge, while in the exact
identity $G''=E[H]-\operatorname{Var}+W_{\mathrm{full}}$ (with $W_{\mathrm{full}}=W_{\mathrm{geo}}+X$ as in
\S\ref{sec:wgeo}) the variance enters with a minus sign. The two faces
differ only in scale (UV versus IR). This unification is stated in the corpus as a
\textbf{structural identification and a conjecture}, not a theorem
(\code{G2-ASSEMBLE-100.md} \S7: ``a structural identification \ldots{} one closure would likely
close both faces --- but that is a conjecture, not a theorem''). This paper preserves that label.

\figscales

\subsection{What this paper is, and is not}

This paper is \textbf{not} a proof of Clay, and it contains no step labelled QED. It is:

\begin{itemize}
\item a precise statement of the reduction, with every hypothesis displayed;
\item a complete inventory of what was actually proved, labelled by type;
\item a statement of the no-gos, each with the precise reason its class fails;
\item the immune map: roughly ninety routes, each with the exact point at which it dies;
\item a specification of the missing tool: five conditions, verbatim;
\item an honest verdict.
\end{itemize}

The reader who wants only the bottom line should read \S\ref{sec:verdict}. The reader who wants
to check the work should follow the file names in \S\ref{sec:theorems}.

\subsection{Roadmap}

\S\ref{sec:halves} states the two halves and the node board with honest percentages.
\S\ref{sec:squeeze} gives the exact spectral squeeze and explains why H1 is necessary but not
sufficient. \S\ref{sec:theorems} states the proven theorems with their hypotheses.
\S\ref{sec:tools} presents the six tools and each one's free-field acid test.
\S\ref{sec:nogos} records the no-gos. \S\ref{sec:species} states the one-species reduction and
the immune map. \S\ref{sec:missing} specifies what the missing tool must satisfy.
\S\ref{sec:verdict} gives the verdict. \S\ref{sec:refs} collects verified references.

% =====================================================================
\section{The two halves: G1 (existence) and G2 (gap), and the node board}
\label{sec:halves}
% =====================================================================

\subsection{G1 --- existence}

G1 asks for a 4D continuum $\mathrm{SU}(N)$ gauge theory satisfying OS0--OS4. The campaign's
decomposition (\code{docs/PROOF-SKELETON.md} \S2) is Table~\ref{tab:g1}.

\begin{table}[t]
\centering
\begin{tabularx}{\textwidth}{@{}l>{\raggedright\arraybackslash}Xc>{\raggedright\arraybackslash}p{3.1cm}@{}}
\toprule
\textbf{Node} & \textbf{Object} & \textbf{Local \%} & \textbf{Label}\\
\midrule
G1.0 & Fixed-volume assembly: L1 definition, L2 theorem, single-block $G''\ge\delta_W$,
  GL3 compactness, GL4 reflection-positivity propagation & $\sim95\%$ (all but L3) & \LT{} mod.\ L3\\
G1.1 & Perturbative/combinatorial kernel: BRST dimension-4 classes reduce to
  $\{\operatorname{Tr}F^2,\operatorname{Tr}F\tilde F\}$; asymptotic freedom absorbs the marginal
  log; $\sum g_j^4=\zeta(4)/(16\beta_0^4)<\infty$ & 100\% & \LT{} closed\\
G1.2 & \textbf{Uniform $Z_a$ $=$ L3}: $\sup_a Z_a^{-1}\|H_a\|_D<\infty$ $=$ L2 $+$ block
  Poincaré $+$ $\Lambda_{\mathrm{UV}}$ tower uniform & $\sim40\%$ & \LM{} open\\
G1.3 & $a\to0$ continuum via OS reconstruction, given L3 & $\sim70\%$ & \LC{} on L3\\
G1.4 & $L\to\infty$, OS2 (continuous O(4)), unique vacuum OS4 (clustering) & 0\% &
  \LM{} open, entangled with G2\\
\bottomrule
\end{tabularx}
\caption{G1 node board: existence.}
\label{tab:g1}
\end{table}

The perturbative kernel G1.1 is a genuine closure: the marginal logarithm is not the obstruction
($\sum g_j^4$ converges), and the only surviving dimension-4 operators are
$\operatorname{Tr}F^2$ and $\operatorname{Tr}F\tilde F$. G1.2 is the sole open UV node, and it
is exactly a \textbf{single-sided force-variance bound} --- the same species as G2's H1. G1.3 is
conditional on G1.2 (OS0/OS1/OS3 are otherwise assembled; OS0 requires L3, OS1 and OS3 are
proved). G1.4 (OS2, infinite volume, uniqueness) is open and, crucially, \textbf{OS4 is
downstream of the gap}: a uniform gap gives clustering for free.

The campaign's honest aggregate: \textbf{G1 is 100\% assembled modulo $\{L3,G2\}$}, with 18
input rows each classified as a theorem, a reduction to L3, the gap, or the same species --- and
no genuinely independent third open item. The master document records a coarser figure of
``G1 $\approx70\%$'' (\code{docs/YANG-MILLS-MASTER.md} \S9); that number is a summary of
``perturbative kernel closed, L3 open,'' not a separate claim. We report both and do not average
them.

\subsection{G2 --- the gap}

G2 asks for $f_1^2>0$ with isolation, equivalently $G''(0)\ge\lambda$ at $k\to0$, uniform in
$(L,a)$. The decomposition (\code{research/unconditional/G2-ASSEMBLE-100.md}) is
Table~\ref{tab:g2}.

\begin{table}[t]
\centering
\begin{tabularx}{\textwidth}{@{}l>{\raggedright\arraybackslash}Xc>{\raggedright\arraybackslash}p{3.1cm}@{}}
\toprule
\textbf{Node} & \textbf{Object} & \textbf{Local \%} & \textbf{Label}\\
\midrule
G2.0 & Identity layer: $G''(k\to0)=E[H|0]-\operatorname{Var}(\nabla S|0)=V_{\mathrm{eff}}''=m_{\mathrm{eff}}^2(k\to0)$;
  exact Ward $E[a\,f]=1$; Stein/IBP $E[H]=E[f^2]$ and $E[\operatorname{Hess}V]=E[\nabla V\otimes\nabla V]\succeq0$ for every measure; $E[-D^2]=E[f^2]$ only for the free field; $G''\operatorname{Var}(\xi)=1$ iff the marginal is Gaussian; the sign sits in the \emph{conditioned} fibre Hessian (cf.\ App.~\ref{app:postaudit}) & 100\% & \LT{} identities\\
G2.1 & Haar geometric term $W_{\mathrm{geo}}=E[\Phi'']-\operatorname{Var}(\Phi')\ge0$
  (Prékopa), every mode including $k\to0$, every $L$ & 100\% & \LT\\
G2.2 & Shallow-field band $\{2g\|F\|\le N/12\}$ closed by $\Phi''\succeq N/12$
  $+$ Brascamp--Lieb & 100\% & \LT\\
G2.3 & High-field cone (WSTC): cone variance mass
  $\le(c\beta\kappa_{\max})^2e^{-c\beta(\kappa^2-\kappa_{\rm typ}^2)N_\square}\to0$,
  $L$-uniform & 100\% & \LT\\
G2.4 & AF-UV deep cone: instability $\le\mu^2e^{-c\mu^4/g^4}\to0$ & 100\% & \LT\\
\textbf{G2.5} & \textbf{Open cell $=$ typical field $\times$ high momentum}:
  the ceiling $\operatorname{Var}(\xi_{k\to0})\le1/\lambda$ (curvature form
  $\operatorname{Var}_\mu(\langle v,\nabla S\rangle|0)\le E_\mu[H_{vv}|0]+W_{\mathrm{full}}-\delta$) &
  \textbf{0\%} & \textbf{\LM{} $=$ H1}\\
\textbf{G2.6} & \textbf{Named species}: negative association / anti-FKG
  $\operatorname{Cov}(G,\Phi'|0)\le0$ & \textbf{0\%} &
  \textbf{\LM{} proven absent from nine classes}\\
\bottomrule
\end{tabularx}
\caption{G2 node board: the gap.}
\label{tab:g2}
\end{table}

The region coverage of the $(\text{field}\times\text{momentum})$ plane is about \textbf{90\%},
closed by the theorems G2.1--G2.4 (Figure~\ref{fig:plane}). The remaining one cell --- typical
field, high momentum --- is a region where all four handles fail \emph{simultaneously}: the WSTC
exponent vanishes when the field is typical, the AF-UV suppression has not begun at
$g\approx1.5$, and the $N/12$ floor is roughly ten times too small. That cell is
\textbf{exactly} $\operatorname{Var}(\xi_{k\to0})\le1/\lambda$, with zero slack. The sign --- the
actual proof --- is \textbf{0\%}.

The G2 assembly is complete in the following precise sense: of 26 input rows, every row is
either (a) a theorem or (b) a reduction to the single species; \textbf{there is no row (c)}, a
genuinely independent open item. The curvature face H1 ($G''\ge\lambda$, the \emph{inverse
zero-momentum susceptibility} $1/\chi$) and the spectral-edge face $H^*$ ($f_1^2\ge r_0$ with
isolation I2) are \emph{not} one object: by the squeeze $f_1^2/\chi\le m_1^2\le M_0/\chi$, H1 is a
susceptibility that \emph{upper}-bounds $m_1^2$ and is necessary but strictly weaker. The
identification $H^*=\text{H1}$ is not merely unproved but is \emph{contradicted} for composite
operators --- the free massive $\mathord{:}\varphi^2\mathord{:}$ has a gap ($m_1=2m$) yet
$G''=1/\chi\ne m_1^2$ --- and every gauge-invariant Yang--Mills operator is composite. The gap is
$H^*$; the rigorous content is the residue--level-gap lemma and the exact Ward identity $E[a\,f]=1$.

\figplane

\subsection{The dimensional correction}

One correction must be recorded because it is easy to get wrong and it changes what
``$G''\ge\lambda$'' means (\code{docs/PROOF-SKELETON.md}, G1-CONTINUUM-SCALING). For the
dimensionless Wilson action, the Hessian satisfies
\begin{equation}
G''_a=a^2\,m_{\mathrm{phys}}^2,\qquad
m_{\mathrm{lat}}=\sqrt{G''_a}=a\,m_{\mathrm{phys}},\qquad
\Delta=\lim m_{\mathrm{lat}}/a .
\label{eq:dim}
\end{equation}
Therefore a claim of the form ``at fixed $a$, $G''_a\ge\lambda$ with $\lambda$ a positive
\emph{constant}'' is dimensionally wrong for the gap: it would give
$m_{\mathrm{phys}}=\lim\sqrt{\lambda}/a=+\infty$, the free-box UV scale, a lattice artifact
rather than a particle. The correct target is a \textbf{ratio} $G''_a/a^2\ge\Delta_0^2$ at the
asymptotically-free scaling rate, i.e.\ $\lambda_a\sim a^2\Delta_0^2$, with $\lambda_a$
transported multiplicatively by $Z_O^2$. The transport itself is free (it is a square); what is
not free is the non-zero seed it multiplies. This distinction pervades the honest reading of the
numerics in \S\ref{sec:tools} and \S\ref{sec:verdict}.

\subsection{The node board, in one table}

\begin{table}[t]
\centering
\begin{tabularx}{\textwidth}{@{}ll>{\raggedright\arraybackslash}p{4.4cm}
  >{\raggedright\arraybackslash}X@{}}
\toprule
\textbf{Half} & \textbf{Object} & \textbf{Local completion} & \textbf{Status}\\
\midrule
G1 & 4D continuum existence & 100\% assembly modulo $\{L3,G2\}$ & \LT{}+\LC, one UV node open\\
G1.2 & UV uniform $Z_a$ $=$ L3 & $\sim40\%$ & \LM\\
G2 & Gap: region coverage & $\sim90\%$ & theorems G2.1--4\\
G2.5--6 & Gap: sign / proof & \textbf{0\%} & \LM{} $=$ H1, the species\\
Channel & $0^{++}$ identification $+$ no dissolution & 100\% & \LC\\
\textbf{Clay} & existence $\wedge$ gap & \textbf{0\%} & \textbf{open}\\
\bottomrule
\end{tabularx}
\caption{The node board in one table (see also Figure~\ref{fig:bars}).}
\label{tab:nodes}
\end{table}

The percentages are \textbf{node-local} --- progress toward each node's own goal --- and are
emphatically \textbf{not} a percentage proof of Clay. Clay is 0\% because the species is
unbuilt.

\figbars
